\SourcePage{736}{739}
\section{Final-state interaction in reactions}

The interaction between particles produced in a reaction may substantially
affect their energy and angular distributions. Naturally, this effect is most
pronounced when their relative speed is small. Such a phenomenon occurs, for
example, in nuclear reactions emitting two or more nucleons, where it is due
to the nuclear forces between the free nucleons\footnote{The results presented
below were obtained by A.~B.~Migdal (1950) and independently by
K.~M.~Watson (1952).}.

Let $\mathbf p_0$ be the momentum of the center of mass of a pair of emitted
nucleons and $\mathbf p$ their relative-motion momentum. Suppose that
$p\ll p_0$, so the relative energy $E=p^2/m$

\SourcePage{737}{740}
($m$ is the nucleon mass) is small compared with the center-of-mass energy
$E_0=p_0^2/4m$. At the same time, suppose $E_0$ is large compared with the
energy $\varepsilon$ of the true or virtual level of the two-nucleon system.
Thus only the relative motion is ``slow''; the nucleons themselves are
``fast.''

The reaction probability is proportional to the squared modulus of the wave
function of the produced particles while they are in the ``reaction zone,''
at separations of order the range $a$ of the nuclear forces (compare the
similar argument for the initial particles in \S~143). Here we seek only the
dependence on the relative motion of one nucleon pair. It is therefore enough
to consider its wave function $\psi_{\mathbf p}(\mathbf r)$, so the
probability of producing a pair with relative momentum in $d^3p$ is
\begin{equation}
 dw_{\mathbf p}=\mathrm{const}\cdot|\psi_{\mathbf p}(a)|^2d^3p.
 \tag{146.1}
\end{equation}

As shown in \S~136, the probability of a scattering transition to a state
with a specified direction of motion requires final-state functions
$\psi_{\mathbf p}^{(-)}$, which contain at infinity a plane wave and only an
incoming spherical wave; they are normalized to a momentum delta-function.
Moreover, $\psi_{\mathbf p}^{(-)}$ follows directly, by complex conjugation
and reversal of $\mathbf p$, from $\psi_{\mathbf p}^{(+)}$, which contains an
outgoing spherical wave and describes mutual scattering of two particles.
This distinction is immaterial in (146.1), so $\psi_{\mathbf p}$ there may be
taken as $\psi_{\mathbf p}^{(+)}$. The problem is thereby reduced to the
resonance scattering of slow particles considered earlier.

Although the true form of $\psi_{\mathbf p}$ for $r\sim a$ is unknown, its
energy dependence can be found by considering it at
$r\gtrsim1/k\gg a$, where $\mathbf k=\mathbf p/\hbar$ and $ak\ll1$, and then
extending its order of magnitude to $r\sim a$\footnote{This procedure is
permissible because for $r\ll1/k$ the energy $E$ may be neglected in the
Schrödinger equation for $\psi_{\mathbf p}$. Its dependence on $E$ in this
region is therefore determined entirely by matching to the function at
$r\sim1/k$.}. In this procedure,

\SourcePage{738}{741}
the spherical wave, containing $1/r$, makes the principal contribution. It is
a sum of partial waves whose amplitudes are the corresponding scattering
amplitudes. For $|\psi_{\mathbf p}(a)|^2$, only the $s$-wave need be retained,
since the amplitudes with $l\ne0$ are relatively small at low energies. From
(133.7), therefore,
\begin{equation}
 \psi_{\mathbf p}\sim\frac1{\kappa+ik}\frac{e^{ikr}}r,
 \tag{146.2}
\end{equation}
where $\kappa=\sqrt{m|\varepsilon|}/\hbar$ and $\varepsilon$ is the energy of
the bound or virtual state of the two-nucleon system\footnote{We mean an
$np$ pair with parallel or antiparallel spins, or an $nn$ pair with
antiparallel spins. For a $pp$ pair, Coulomb repulsion complicates the
situation; that case must be treated by the theory of \S~138.}. Substitution
in (146.1) gives
\begin{equation}
 dw_{\mathbf p}=\mathrm{const}\cdot\frac{d^3p}{E+|\varepsilon|}.
 \tag{146.3}
\end{equation}

Thus the distribution in momentum direction, in the two-nucleon
center-of-mass system, is isotropic. The relative-energy distribution is
\begin{equation}
 dw_E=\mathrm{const}\cdot\frac{\sqrt E\,dE}{E+|\varepsilon|}.
 \tag{146.4}
\end{equation}
The nucleon interaction therefore produces a maximum at small $E$, with
$E\sim|\varepsilon|$\footnote{Strictly, the constant coefficients in
(146.3), (146.4) may also depend on $E$ through the other parts of the wave
function of the complete reaction-product system. This dependence is weak:
as a function of $E$, the coefficient changes appreciably only over the full
energy interval of order $E_0$ available to the nucleon pair. In the region
$E\ll E_0$, it may therefore be neglected beside the strong dependence in
(146.4).}.

Small relative momentum, $p\ll p_0$, corresponds in the laboratory system to
small angles $\theta$ between the two nucleon momenta. The maximum in the
energy distribution therefore appears as an angular correlation between the
emission directions, with enhanced probability at small $\theta$.

\SourcePage{739}{742}
Let $\mathbf p_1$, $\mathbf p_2$ be the laboratory momenta. Then
\[
 \mathbf p_0=\mathbf p_1+\mathbf p_2,\qquad
 \mathbf p=\frac12(\mathbf p_2-\mathbf p_1),
\]
since the reduced mass of two identical particles is $m/2$. Taking the vector
product of these equalities gives
$\mathbf p_0\times\mathbf p=\mathbf p_1\times\mathbf p_2$. For $p\ll p_0$,
\[
 p_0p_\perp=p_1p_2\sin\theta\approx\frac{p_0^2}{4}\theta,
\]
or $\theta=4p_\perp/p_0$, where $p_\perp$ is the component of $\mathbf p$
transverse to $\mathbf p_0$. Rewrite (146.3) as
\[
 dw_{\mathbf p}=\mathrm{const}\cdot
 \frac{2\pi p_\perp\,dp_\perp\,dp_\parallel}
 {(1/m)(p_\perp^2+p_\parallel^2)+|\varepsilon|}.
\]
Integration over $dp_\parallel$ gives the angular distribution. Because the
integral converges rapidly, its limits may be extended from $-\infty$ to
$\infty$, giving
\begin{equation}
 dw_\theta=\mathrm{const}\cdot
 \frac{\theta\,d\theta}{\sqrt{\theta^2+4|\varepsilon|/E_0}}.
 \tag{146.5}
\end{equation}
Per unit solid angle $d\Omega\approx2\pi\theta\,d\theta$, the angular
distribution has a maximum at
$\theta\sim\sqrt{|\varepsilon|/E_0}$.
